\documentclass[11pt]{article}
\usepackage[margin=1in]{geometry}
\usepackage{parskip}
\usepackage[colorlinks=true,allcolors=blue]{hyperref}

\begin{document}

\pagestyle{empty}

\noindent Hikaru Wakaura and Taiki Tanimae\\
QIRI (Quantum Integrated Research Institute Inc.)\\
Tokyo 107-0061, Japan\\
\href{mailto:h.wakaura@deeptell.jp}{h.wakaura@deeptell.jp} \quad ORCID: 0000-0001-8381-8323

\vspace{1em}
\noindent July 9, 2026

\vspace{1em}
\noindent To the Editors, \emph{PRX Quantum}

\vspace{1em}

\noindent \textbf{Re: Submission of ``Chaotic states outperform quantum many-body
scars as codes under generic measurement: a refutation of scar-enhanced
information protection''}

\vspace{0.5em}

On July~9, 2026, we submit for your consideration the manuscript above (research
article, two authors) together with Supplemental Material (one PDF: eight
sections, three figures, one table) and an openly available simulation package
reproducing every result.

\vspace{0.5em}

% Paragraph 1: what we found
Quantum many-body scars---rare non-thermal, weakly entangled eigenstates that
resist thermalization---are widely regarded as promising long-lived carriers of
quantum information. We put that intuition to a quantitative test and overturn
it. Encoding a logical qubit (and, more generally, the largest code a scar
subspace admits) in the scar states of a monitored circuit and measuring the
reference-qubit coherent information, we find that generic chaotic (thermal)
states protect the encoded information \emph{better} than scars under every
measurement we examined, at every system size, in two independent models. A
striking part of the story is methodological: a comparison against a single,
seemingly natural thermal reference state reports the opposite conclusion, and
only ensemble averaging reveals that this reference is an atypical outlier.

% Paragraph 2: why it matters (message + surprise)
The message is that the resource protecting quantum information under measurement
is a code's spread in the measurement basis---scrambling---not its athermality;
weak ergodicity breaking is a liability rather than an asset. The mechanism is
sharp and, we believe, memorable: the emergent symmetry that makes scars special
is exactly what dooms them, because its generator coherently couples the encoded
states and its Casimir carries an irreducible variance that a measurement can
resolve. We turn this into a falsifiable prediction and confirm it: in a second,
exactly-solvable model whose scars close the algebra \emph{exactly}, the same
Casimir measurement instead protects the scar perfectly, isolating the variance
as the single control parameter. Beyond the physics, the ensemble-baseline
reversal is a cautionary lesson of broad utility for the many groups now
benchmarking codes and channels against individual eigenstates.

% Paragraph 3: why this journal
We believe \emph{PRX Quantum} is the natural home for this work. It operationally
realizes, directly in the monitored-circuit setting, the correspondence between
eigenstate thermalization and approximate quantum error correction; it removes
weak-ergodicity-breaking subspaces from the design space of measurement-robust
quantum memories; and it delivers a clean, reproducible, mechanism-level result
of the conceptual significance and rigor that define the journal's scope. All
code and data are openly available for one-command reproduction.

\vspace{0.5em}

\noindent We suggest the following referees, all leading experts in
measurement-induced dynamics, quantum error correction, or quantum many-body
scars (institutional contact addresses are publicly listed):

\begin{itemize}\itemsep2pt
  \item Prof.\ Matthew P.~A.\ Fisher, University of California, Santa Barbara ---
        measurement-induced transitions
  \item Prof.\ Ehud Altman, University of California, Berkeley ---
        error correction and monitored circuits
  \item Dr.\ Michael J.~Gullans, NIST / University of Maryland ---
        purification transitions
  \item Prof.\ Vedika Khemani, Stanford University --- random and monitored
        circuits
  \item Prof.\ Adam Nahum, \'Ecole Normale Sup\'erieure, Paris --- entanglement
        dynamics
  \item Prof.\ Zlatko Papi\'c, University of Leeds --- quantum many-body scars
  \item Prof.\ Maksym Serbyn, IST Austria --- quantum many-body scars
  \item Dr.\ Sanjay Moudgalya, Technical University of Munich --- exact scars and
        spectrum-generating algebras
\end{itemize}

\noindent We do not request the exclusion of any referee.

\vspace{1em}

\noindent On behalf of both authors,\\
Hikaru Wakaura\\
QIRI (Quantum Integrated Research Institute Inc.), Tokyo 107-0061, Japan\\
\href{mailto:h.wakaura@deeptell.jp}{h.wakaura@deeptell.jp}

\end{document}
